\section{Géométries amplituédriques de rang un et compatibilité des résidus}
\label{sec:k-poligono-amplituhedral}

Le registre dodécaphasique permet de réunir explicitement des données positives,
un domaine grassmannien, une image cinématique, un recouvrement, un degré sur les cellules
et une forme canonique. Le premier résultat concerne l’amplituèdre
bidimensionnel de paramètres \((n,k,m)=(13,1,2)\) ; la construction est ensuite
prolongée aux polytopes cycliques de rang un et de dimension supérieure.
Il s’agit d’une réalisation géométrique ultérieure du registre produit
par APP--TRIT--TPK. L’opération supplémentaire déclarée est l’élévation aux
puissances d’une partition rationnelle ordonnée. Ce choix fixe une
réalisation particulière ; les correspondances physiques ultérieures conservent
leurs lecteurs et leurs conditions propres.

La provenance du registre conserve la composition
\[
 x_{\mathrm{term}}\longmapsto\mathcal R_{12}(x_{\mathrm{term}})
 \longmapsto(b^{90},b^{120},Q_{\mathrm{TPK}})
 \longmapsto U_{\mathrm{sgn}}\longmapsto K.
\]
APP conserve les deux feuillets et leurs résidus et quotients ; TRIT oriente le régime de chaque transition ; TPK transporte et actualise l'état avec retenue, frontière, route et mémoire. La représentation \(K\) se lit dans le même système d'états enrichis qui engendre conjointement le continu. Son extraction signée et son inversion de Hadamard produisent
\begin{equation}
 K=(234,543,140,729,659,824,621,58,914,794,146,601),
 \qquad S_K=\sum_{j=1}^{12}K_j=6263.
 \label{eq:polygon-K}
\end{equation}
La preuve géométrique utilise la positivité de ces douze coordonnées, leur
ordre et leur total. Ses coefficients proviennent de cette extraction et non de
la sélection rétrospective d’une valeur physique.

\subsection{Partition rationnelle et relèvement positif}

Nous définissons treize points rationnels de partition :
\begin{equation}
 t_0=0,\qquad t_i=\frac{K_1+\cdots+K_i}{S_K}\quad(1\leq i\leq12),
 \qquad Z_i=(1,t_i,t_i^2)^\mathsf T,\qquad Z=(Z_0\ \cdots\ Z_{12}).
 \label{eq:polygon-Z}
\end{equation}
Leurs numérateurs, de dénominateur commun \(6263\), sont
\[
 (0,234,777,917,1646,2305,3129,3750,3808,4722,5516,5662,6263).
\]
Ainsi, \(0=t_0<t_1<\cdots<t_{12}=1\). La partition étiquetée et le total conservent le registre au moyen de l'inverse exact
\begin{equation}
 K_i=S_K(t_i-t_{i-1}).
 \label{eq:polygon-recover-K}
\end{equation}
La partition conserve les proportions ; retenir \(S_K\) conserve également
l’échelle entière.

\begin{proposition}[Données cinématiques strictement positives]
\label{prop:polygon-vandermonde}
La matrice \(Z\) a rang trois et tous ses mineurs ordonnés de taille
trois sont positifs. Pour \(0\leq a<b<c\leq12\),
\begin{equation}
 \langle abc\rangle:=\det(Z_a,Z_b,Z_c)
 =(t_b-t_a)(t_c-t_a)(t_c-t_b)>0.
 \label{eq:polygon-minors}
\end{equation}
\end{proposition}
\begin{proof}
On soustrait la première colonne des deux dernières et on développe selon la
première ligne. Chaque différence \(t_j^2-t_a^2\) se factorise en
\((t_j-t_a)(t_j+t_a)\), ce qui donne le produit indiqué. Les sommes partielles
strictement croissantes prouvent sa positivité. Chacun de ces
mineurs atteste le rang trois.
\end{proof}

\subsection{Image, fibres et cellules de degré un}

La grassmannienne \(\operatorname{Gr}_{\geq0}(1,13)\) est l'espace projectif des vecteurs \(c=(c_0,\ldots,c_{12})\) non négatifs et non nuls, modulo une échelle positive. La coordinatisation \(\sum c_j=1\) l'identifie au simplexe fermé de dimension douze. On définit
\begin{equation}
 \Phi_Z:\operatorname{Gr}_{\geq0}(1,13)\longrightarrow\mathbb P^2,
 \qquad[c]\longmapsto\left[\sum_{j=0}^{12}c_jZ_j\right].
 \label{eq:polygon-map}
\end{equation}
Dans la coordinatisation \(Y_0=1\), soient \(v_j=(t_j,t_j^2)\) et \(P_K=\operatorname{conv}(v_0,\ldots,v_{12})\).

\begin{theorem}[Amplituèdre polygonal du registre]
\label{thm:polygon-amplituhedron}
L’image fermée de \(\Phi_Z\) est exactement le polygone convexe
\(P_K=\mathcal A_{13,1,2}(Z)\), de dimension deux et à treize sommets. L’image
de la grassmannienne strictement positive est l’intérieur de \(P_K\).
Les onze cellules
\begin{equation}
 \mathcal S_i=\{[c]:c_0,c_i,c_{i+1}>0,\
 c_j=0\text{ si }j\notin\{0,i,i+1\}\},\qquad 1\leq i\leq11,
 \label{eq:polygon-cells}
\end{equation}
s'appliquent bijectivement, avec un degré un et une orientation positive, sur les intérieurs de \(T_i=[v_0,v_i,v_{i+1}]\). Leurs adhérences recouvrent \(P_K\) et leurs intérieurs sont deux à deux disjoints.
\end{theorem}
\begin{proof}
Dans la normalisation \(\sum c_j=1\), l’application est
\(c\mapsto\sum c_jv_j\), dont l’image est l’enveloppe convexe. Les mineurs
de \eqref{eq:polygon-minors} montrent que chaque triplet ordonné a une
orientation positive. Chaque côté consécutif laisse tous les autres sommets
strictement à sa gauche ; il en va de même du côté final allant de \(v_{12}\)
à \(v_0\). Tous les points sont donc des sommets extrémaux d’un polygone
convexe.

Toute combinaison avec \(c_j>0\) est intérieure, car pour toute droite
d’appui, l’un des sommets se trouve strictement dans le demi-plan
intérieur. Réciproquement, soient \(y\) intérieur et
\(b=13^{-1}\sum v_j\). Pour \(\varepsilon>0\) suffisamment petit,
\(y'=(y-\varepsilon b)/(1-\varepsilon)\) demeure dans le polygone.
En écrivant \(y'=\sum a_jv_j\), \(a_j\geq0\), \(\sum a_j=1\), on obtient
\(y=\sum((1-\varepsilon)a_j+\varepsilon/13)v_j\), avec tous les
coefficients positifs.

Les diagonales issues de \(v_0\) divisent le polygone en onze triangles. Sur chaque support triple, les coordonnées barycentriques sont uniques, car \(\langle0,i,i+1\rangle>0\) ; leur inverse explicite apparaît dans \eqref{eq:polygon-barycentric}. Chaque support définit une cellule positroïdale de rang un de coordonnées \([1:a:b]\), \(a,b>0\), à ces positions. Un réseau orienté planaire comportant une source et trois trajectoires de poids \(1,a,b\) vers les trois destinations réalise ces coordonnées par mesure de frontière. Le déterminant positif prouve l'orientation et le degré un.
\end{proof}

La fibre générique de l’application complète est de dimension dix :
\(\sum c_j=1\), \(\sum c_jt_j=x\) et \(\sum c_jt_j^2=y\) sont trois
équations indépendantes en treize coordonnées. Pour \(y\) intérieur, la
fibre contient un point strictement positif, de sorte que son intersection
avec le simplexe est de dimension \(13-3=10\). Le degré un concerne les
cellules triangulaires ; la forme bidimensionnelle s’obtient au moyen de ces
cellules.

\subsection{Forme canonique et résidus orientés}

Pour \(Y=(1,x,y)^\mathsf T\), nous définissons
\begin{equation}
 \ell_{ab}(Y)=\det(Z_a,Z_b,Y)
 =(t_b-t_a)\{y-(t_a+t_b)x+t_at_b\}.
 \label{eq:polygon-lines}
\end{equation}
L’expression vaut également lorsque \(b<a\). Pour \(a<b<c\), soient
\(D_{abc}=\langle abc\rangle\) et
\begin{equation}
 \lambda_a=\frac{\ell_{bc}}{D_{abc}},\qquad
 \lambda_b=\frac{\ell_{ca}}{D_{abc}},\qquad
 \lambda_c=\frac{\ell_{ab}}{D_{abc}},\qquad
 \lambda_a+\lambda_b+\lambda_c=1.
 \label{eq:polygon-barycentric}
\end{equation}
L’orientation affine est \(dx\wedge dy\). On fixe la convention de frontière
\(\operatorname{Res}^{\partial}(\eta\wedge du/u)=\eta|_{u=0}\) :
la différentielle normale occupe la dernière position. En dimension deux, le
résidu de Poincaré qui place \(du/u\) en premier a le signe opposé.
Les annulations suivantes utilisent uniformément la convention déclarée.

\begin{theorem}[Forme rationnelle et annulation des diagonales]
\label{thm:polygon-form}
La forme canonique du triangle orienté \([v_a,v_b,v_c]\) est
\begin{equation}
 \Omega_{abc}=
 \frac{D_{abc}^2\,dx\wedge dy}
 {\ell_{ab}(Y)\ell_{bc}(Y)\ell_{ca}(Y)}.
 \label{eq:polygon-triangle-form}
\end{equation}
La forme canonique du polygone est
\begin{equation}
 \boxed{\displaystyle
 \Omega_{P_K}(x,y)=
 \sum_{i=1}^{11}
 \frac{t_it_{i+1}(t_{i+1}-t_i)\,dx\wedge dy}
 {(y-t_ix)\{y-(t_i+t_{i+1})x+t_it_{i+1}\}(t_{i+1}x-y)}.}
 \label{eq:polygon-full-form}
\end{equation}
Les pôles des diagonales intérieures s’annulent. Les seuls diviseurs
polaires sont les treize droites des côtés extérieurs. Chaque côté étant paramétré
de son sommet initial à son sommet final par \(s\in[0,1]\), son résidu de
frontière est \(ds/[s(1-s)]\).
\end{theorem}
\begin{proof}
La forme du simplexe barycentrique est \(d\lambda_b\wedge d\lambda_c/(\lambda_a\lambda_b\lambda_c)\). Le jacobien satisfait \(dx\wedge dy=D_{abc}\,d\lambda_b\wedge d\lambda_c\). Substituer \eqref{eq:polygon-barycentric} donne \eqref{eq:polygon-triangle-form}. Pour \(a=0\), simplifier les trois facteurs scalaires des droites donne chaque terme de \eqref{eq:polygon-full-form}.

Sur le côté \(ab\), posons \(\lambda_c=u\),
\(\lambda_b=(1-u)s\), \(\lambda_a=(1-u)(1-s)\). Alors
\[
 \Omega_{abc}=\frac{ds}{s(1-s)}\wedge\frac{du}{u(1-u)}.
\]
Le résidu est la forme de l’intervalle orienté. Deux triangles contigus
induisent des orientations contraires sur leur diagonale et produisent des résidus
opposés. Tous les pôles étant simples, annuler le résidu élimine le
pôle de la diagonale comme diviseur. Chaque côté extérieur apparaît une fois.
La formule homogène est de degré zéro comme forme projective et n’a pas
d’autres diviseurs polaires. Si deux formes rationnelles projectives avaient les
mêmes résidus et aucun autre pôle, leur différence serait une deux-forme
holomorphe sur \(\mathbb P^2\), qui est nécessairement nulle. La
caractérisation et l’unicité sont ainsi prouvées.
\end{proof}

On emploie la notion conventionnelle de géométrie positive au moyen de pôles logarithmiques et de résidus canoniques de frontière, développée par N.~Arkani-Hamed, Y.~Bai et T.~Lam, \emph{Positive Geometries and Canonical Forms}, JHEP 11 (2017), 039, \url{https://arxiv.org/abs/1703.04541}, sections 2, 3, 5.2, 6.1 et 7.2. La sélection des coefficients du présent polygone précède cette reconnaissance et provient de \(K\).

\subsection{Raffinement par retenue sur la géométrie fixe}

Dans un triangle \(T=[A,B,C]\) d’orientation positive, nous écrivons
\begin{equation}
 Y(u,r)=(1-u)A+u\{(1-r)B+rC\},\qquad 0<u,r<1.
 \label{eq:polygon-blowdown}
\end{equation}
Le sommet \(A\) correspond à \(u=0\) et le côté opposé à \(u=1\).
La forme devient
\begin{equation}
 \Omega_T=\frac{du}{u(1-u)}\wedge\frac{dr}{r(1-r)}.
 \label{eq:polygon-product-form}
\end{equation}
Soit \(B_0\geq2\) une base entière déclarée, distincte du sommet \(B\).
La coordonnée \(r\) admet la mise à jour par résidu et retenue
\begin{equation}
 c=\lfloor B_0r\rfloor,\qquad r'=B_0r-c,\qquad
 r=\frac{c+r'}{B_0}.
 \label{eq:polygon-carry}
\end{equation}
Les extrémités conservent leurs marques de frontière. Le chiffre \(c\) demeure comme mémoire et permet d'inverser l'actualisation. Il convient de préciser les coordinatisations des deux transports affines. Le zoom de la sous-cellule agit sur \(r\in[c/B_0,(c+1)/B_0]\) au moyen de \(r'=B_0r-c\). En revanche, le transport de la paire relevée \(L_c(x,y)=(B_0x-c,B_0y+c)\), avec \(x+y=M\), change la masse totale en \(B_0M\) et exprime la coordonnée
\[
 \frac{B_0x-c}{B_0M}=\frac{x}{M}-\frac{c}{B_0M}.
\]
Ce sont des applications de domaines et de normalisations différents. Toutes deux conservent la retenue dans l'état enrichi et permettent une récupération avec leurs données d'échelle. La démonstration suivante utilise le zoom de sous-cellule et la partition géométrique qu'il identifie.

\begin{theorem}[Invariance exacte sous subdivision avec mémoire]
\label{thm:polygon-carry}
Soient \(s_j=j/B_0\) et \(W_j=(1-s_j)B+s_jC\).
Les triangles \(T_j=[A,W_j,W_{j+1}]\), \(0\leq j<B_0\), subdivisent
\(T\) et satisfont
\begin{equation}
 \sum_{j=0}^{B_0-1}\Omega_{T_j}=\Omega_T.
 \label{eq:polygon-carry-additivity}
\end{equation}
L'égalité reste valable à toute profondeur finie et pour toute partition rationnelle ordonnée du côté. Subdiviser les onze triangles de \(P_K\) laisse \(\Omega_{P_K}\) exactement invariante.
\end{theorem}
\begin{proof}
Les régions \(r\in[s_j,s_{j+1}]\) couvrent \(T\) avec des intérieurs disjoints.
Dans les coordonnées communes,
\[
 \Omega_{T_j}=\frac{du}{u(1-u)}\wedge
 \frac{(s_{j+1}-s_j)\,dr}{(r-s_j)(s_{j+1}-r)}.
\]
L’identité rationnelle
\[
 \frac{s_{j+1}-s_j}{(r-s_j)(s_{j+1}-r)}
 =\frac1{r-s_j}+\frac1{s_{j+1}-r}
\]
se télescope : chaque terme intérieur annule celui du sous-intervalle contigu.
Il reste \(1/r+1/(1-r)=1/[r(1-r)]\). Cela prouve l’égalité des formes
rationnelles ; son itération fournit toute profondeur finie.
La preuve utilise seulement \(s_0=0<s_1<\cdots<s_N=1\), de sorte qu’elle couvre
les partitions non uniformes. Le changement
\(r'=(r-s_j)/(s_{j+1}-s_j)\) restitue localement la forme unitaire.
Le mot des retenues identifie la sous-cellule et son inverse récupère la
coordonnée antérieure. Les nouveaux pôles intérieurs s’annulent au moyen de
la même identité télescopique.
\end{proof}

La somme précédente réunit la partition géométrique complète. Lorsqu’un lecteur
de survie sélectionne seulement certaines histoires, son image est l’union
des sous-cellules sélectionnées et sa forme est la somme correspondante ;
elle coïncide avec la forme du polygone complet précisément lorsque cette union
le recouvre avec une multiplicité orientée de un. Le critère de
recouvrement pour une sélection concrète d’histoires s’exprime donc par
une condition vérifiable sur sa chaîne de cellules.

Ce raffinement conserve le polygone fixe. Insérer un nouveau point \((t,t^2)\) de la parabole entre deux paramètres consécutifs \(a<t<b\) produit une autre géométrie : la hauteur de la corde moins \(t^2\) est \((t-a)(b-t)>0\), et le nouveau point se trouve à l'extérieur du polygone précédent. L'invariance démontrée correspond à la subdivision rectiligne des cellules, dont la frontière extérieure reste fixe.

\subsection{Changements de triangulation et reparamétrage positif}

Pour quatre sommets ordonnés \(A,B,C,D\), on a
\begin{equation}
 \Omega_{ABC}+\Omega_{ACD}=\Omega_{ABD}+\Omega_{BCD}.
 \label{eq:polygon-flip}
\end{equation}
Les deux membres ont les mêmes résidus extérieurs et annulent leur diagonale intérieure ; leur égalité découle de l'unicité démontrée. Elle se vérifie aussi en multipliant par leurs dénominateurs linéaires. Tout changement entre triangulations d'un polygone convexe se décompose en ces changements de diagonale, si bien que la forme est indépendante de la triangulation. Une identification avec des coordonnées cluster doit en outre conserver la variété et ses relations de changement.

Si \(D=\operatorname{diag}(d_0,\ldots,d_{12})\), \(d_j>0\), alors
\begin{equation}
 \mathcal A_{13,1,2}(ZD)=\mathcal A_{13,1,2}(Z).
 \label{eq:polygon-rescale}
\end{equation}
L'automorphisme \([c]\mapsto[Dc]\) reparamétrise l'application cinématique. Dans la coordinatisation affine, \(c_j\) devient \(c_jd_j/\sum_\ell c_\ell d_\ell\), qui parcourt de nouveau le simplexe. L'image et sa forme canonique demeurent identiques. Le déplacement d'un historique paramétré et le changement de l'ensemble géométrique complet sont ainsi des opérations distinctes.

\subsection{Extension de rang un aux dimensions supérieures}

Pour \(2\leq m\leq12\), nous définissons
\begin{equation}
 Z_i^{(m)}=(1,t_i,t_i^2,\ldots,t_i^m)^\mathsf T,\qquad
 P_K^{(m)}=\operatorname{conv}
 \{(t_i,t_i^2,\ldots,t_i^m):0\leq i\leq12\}.
 \label{eq:polygon-higher-lift}
\end{equation}
Les mineurs maximaux ordonnés sont les produits de Vandermonde
\(\prod_{a<b}(t_{i_b}-t_{i_a})>0\). La matrice a rang \(m+1\).
L’image de \(\operatorname{Gr}_{\geq0}(1,13)\) est le polytope cyclique
\(P_K^{(m)}=\mathcal A_{13,1,m}(Z^{(m)})\), de dimension \(m\) ; sa fibre
générique a dimension \(12-m\). On utilise ici l’extension mathématique
de la définition amplituédrique à tout \(m\) ; les applications
habituelles aux amplitudes imposent leurs propres valeurs de \(m\).

La triangulation est déterminée par l’ordre fixe des étiquettes. On choisit
le plus petit sommet d’un polytope, on triangule récursivement toutes ses
facettes qui ne le contiennent pas et on prend les cônes depuis ce sommet. La
récurrence commence par les points. Une droite allant du sommet initial
à un point intérieur coupe la frontière opposée dans une facette ; cela
prouve le recouvrement. L’unicité du point d’intersection et la
compatibilité des restrictions induites par le même ordre prouvent
que les intérieurs sont disjoints et que les intersections sont des faces
communes. À chaque étape, la dimension diminue, de sorte que le procédé
s’achève. Il s’agit d’une triangulation par traction déterminée par les
étiquettes.

Pour un simplexe orienté \(T=[p_0,\ldots,p_m]\) de cette triangulation, soient
\[
 A_T=(p_1-p_0\ \cdots\ p_m-p_0),\qquad
 (\lambda_1,\ldots,\lambda_m)^\mathsf T=A_T^{-1}(x-p_0),
 \qquad\lambda_0=1-\sum_{j=1}^m\lambda_j.
\]
L’orientation de ses sommets est choisie de sorte que \(\det A_T>0\).
Sa forme canonique est
\begin{equation}
 \Omega_T(x)=
 \frac{dx_1\wedge\cdots\wedge dx_m}
 {\det(A_T)\prod_{j=0}^m\lambda_j(x)},\qquad
 \Omega_{P_K^{(m)}}=\sum_{T}\Omega_T.
 \label{eq:polygon-higher-form}
\end{equation}
Le changement barycentrique prouve la formule et le fait que la restriction cinématique sur le support de \(T\) est de degré un. Les résidus sont les formes des faces avec les signes induits par l'orientation. Sur une face intérieure, les deux orientations sont opposées et le résidu total s'annule. Sur une face extérieure reste la somme des formes de la triangulation de cette face. La récurrence sur la dimension démontre que cette somme est sa forme canonique. Répéter le résidu le long d'un drapeau de faces atteint les sommets, où la normalisation est \(\pm1\). On obtient ainsi la correspondance des résidus dans toutes les codimensions pour cette collection de rang un.

La même démonstration couvre toute subdivision simpliciale compatible du polytope fixé : ses faces intérieures s'annulent et les faces extérieures conservent les résidus. L'absence de formes holomorphes de degré maximal sur \(\mathbb P^m\) assure l'unicité de la forme rationnelle. En particulier, la stabilité par raffinement, la couverture et les résidus sont établis dans cette collection de dimensions \(2,\ldots,12\). Les degrés grassmanniens \(k>1\) conservent leurs applications et leurs obligations démonstratives spécifiques.
